% !TeX program = XeLaTeX
% !TeX root = ../pujavidhanam.tex
% --meta--
% source: श्रीगङ्गाकल्पोक्तपूजाविधिः, ed. Mittur Srinivasa Bhatta (Brahmashri Mittur Purohita
%   Timmayya Bhatta Samprathishthana, Idkidu, 2015), pp. 1-40; from the family's handwritten
%   paddhati. The pūjā done on return from Kāśī-yātrā (kāśī-samārādhana), on Gaṅgāṣṭamī, or
%   when a new well or borewell first gives water. Rows marked for review: review/ganga/.
% --end-meta--
\chapt{श्री-गङ्गा-कल्पोक्त-पूजा}

\begin{center}
॥श्रीगणेशाय नमः॥ ॥श्रीगङ्गाविश्वेश्वराय नमः॥
\end{center}

अथ कर्ता सपत्नीकः मङ्गलस्नानादि नित्यकर्म कृत्वा भद्रासने उपविश्य आचम्य (प्राणानायम्य)
पवित्रपाणिः। (गुर्वाद्यभिवन्दनं कृत्वा) ॐ ह्रीं परमात्माद्यशेषगुरुपारम्पर्यक्रमेण
स्वगुरुपादाम्बुजं यावत्तावत्प्रणौमि इत्यादि कलशपूजान्तं कृत्वा `सुमुखश्चैकदन्तश्च' इत्याद्युक्त्वा
सङ्कल्पं कुर्यात्।

\dnsub{सङ्कल्पः}

ॐ स्वस्ति। \ldots\ शुभतिथौ। काशी-गया-प्रयागादि-तीर्थयात्रायाः प्रतिनिवृत्तस्य \blank{}
गोत्रोत्पन्नस्य \blank{} नक्षत्रे \blank{} राशौ जातस्य \blank{} (शर्मणः) मम भार्यायाः \blank{}
अस्माकं सहकुटुम्बानां क्षेम-स्थैर्य-धैर्य-वीर्य-विजय-शुभ-अभय-आनन्द-आरोग्य-ऐश्वर्याभिवृद्ध्यर्थं
धर्म-अर्थ-काम-मोक्षाख्य-चतुर्विधफलपुरुषार्थसिद्ध्यर्थम् आदित्यादीनां नवानां ग्रहाणां
शुभैकादशस्थानफलावाप्त्यर्थं मम पितृमातृवंशजातानां दिवि भुवि तिर्यङ्नरकादिवासिनाम्
अन्येषां च अस्मत्पूर्वजानाम् उत्तमलोकावाप्त्यर्थं समस्तपूर्वजोद्धारार्थं
नान्दीमुखाश्रुमुखादि-पितृप्रीत्यर्थम् ऊष्मपो यमोऽङ्गिरस्वानग्निष्वात्तोऽग्निः कव्यवाहनादयो
ये पितरस्तेषामविच्छिन्नप्रसादसिद्ध्यर्थम् असृग्भागि-लेपभागि-पिण्डभागि-विप्रोच्छिष्टभागि-प्रभृतीनाम्
अक्षय्यतृप्त्यर्थं नानापुण्यनदीस्नानपूजाफलसिद्ध्यर्थं मम जन्मप्रभृति अद्य यावत्
ज्ञाताज्ञात-कामाकाम-सकृदसकृत्कृत-कायिक-वाचिक-मानसिक-स्पृष्टास्पृष्ट-पीतापीत-भुक्ताभुक्त-सकलपातकोपपातक-लघुपातक-सङ्करीकरण-मलिनीकरणापात्रीकरण-जातिभ्रंशकर-प्रकीर्णकादीनां
स्वतः परप्रेरणया वा सम्भावितानां सकलविधपापानां परिहारार्थं यमुना-सरस्वती-चन्द्रभागाख्य-महानदीस्नान-कल्पोक्तफलसिद्ध्यर्थं
वेत्रवती-मलापहा-फल्गुनी-कृष्णाख्य-पवित्रनदीक्षेत्रयात्राफलसिद्ध्यर्थं
शुतुद्री-गोदावरी-मरुद्वृधाख्य-पुण्यनदीतीर्थमज्जन-शतफलसिद्ध्यर्थं
भीमरथी-सीता-चर्मण्वती-नर्मदाख्य-विशिष्टनदीमाहात्म्यश्रवणफलावाप्त्यर्थं
पयोष्णी-विपाशा-ताम्रपर्णेत्यादि-तत्तन्नदीकल्पोक्त-तत्तत्फलसिद्ध्यर्थं
कावेरी-करतोया-नेत्रावत्यादि-नदीविशेष-तीर्थविशेष-स्नान-तर्पण-फलावाप्त्यर्थं
काश्यां विश्वनाथदर्शनस्य पूर्णफलावाप्त्यर्थं गयायां पिण्डदानस्य साकल्यार्थं
प्रयागे त्रिवेणीसङ्गमस्नानस्य समग्रतावाप्त्यर्थं गङ्गाप्रीत्यर्थं गङ्गाकल्पोक्तफलसिद्ध्यर्थं
निर्वर्तितयात्रोऽहम् अस्मिन् गङ्गातीर्थोदके गङ्गाकल्पोक्तपूजां करिष्ये॥

{\small [नूतनतया प्रविधारितजलप्रणाल्यां (बोरवेल्): `\ldots\ जलस्य शिवगङ्गात्वसिद्ध्यर्थम् अस्मिन्
नवोदके गङ्गापूजां करिष्ये'। नवोत्खातकूपात् कुल्याया वा गृहीतजले: `नवोत्खातकुल्याया जलस्य
शिवगङ्गात्वसिद्ध्यर्थम् अस्मिन् नवोदके गङ्गापूजां करिष्ये'।]\par}

तदादौ गणपतिपूजां करिष्ये। इति सङ्कल्प्य, तत्कृत्वा। क्षेत्रपालं वास्तोष्पतिं च सम्पूज्य।
सुवर्णपुष्पं दत्वा। ब्राह्मणार्चनं कृत्वा (पञ्चवाक्यैः पुण्याहं वाचयित्वा; तत्र `श्रीगङ्गा
(विश्वेश्वरः) प्रीयताम्' इति विशेषः)। `आगमार्थं तु' इत्यादि घण्टानादं कृत्वा कलशपूजा,
शङ्खपूजा, दीपपूजा, आत्मपूजा इत्यादि कृत्वा पीठपूजां कृत्वा [विश्वेश्वरस्य नवशक्तिपूजां कृत्वा
(वामायै नमः, ज्येष्ठायै नमः\ldots) विश्वेश्वरम् आवाह्य सम्पूज्य]।

`ॐ मही द्यौः पृथिवी' इत्यादि कलशं संस्थाप्य, वरुणं च आवाह्य सम्पूज्य, `ॐ आपो हि ष्ठा',
`हिरण्यवर्णाम्' इत्यादि मन्त्रैः प्रोक्ष्य, अथ

\sect{नवशक्ति-पूजा}

श्रीगङ्गायै नमः, नवशक्तिपूजां करिष्ये।

\begin{center}
ॐ जयायै नमः। ॐ विजयायै नमः। ॐ अजितायै नमः।\\
ॐ अपराजितायै नमः। ॐ नित्यायै नमः। ॐ विलासिन्यै नमः।\\
ॐ दोग्ध्र्यै नमः। ॐ अघोरायै नमः। ॐ मङ्गलायै नमः।\\
\textbf{नवशक्तिपूजां समर्पयामि।}
\end{center}

ॐ नमो भगवत्यै सर्वभूतात्मिकायै सकलगुणशक्तियुक्तायै त्रिपथगामिन्यै योगपद्मपीठात्मिकायै
गङ्गायै नमः, सुवर्णपीठं कल्पयामि।

\twolineshloka*
{स्वात्मसंस्थामजां शुद्धां त्वामद्य परमेश्वरि}
{अरण्यामिव हव्याशं (मूर्त्याम्) अप्सु आवाहयाम्यहम्}

आवाहयेत्। ॐ भूः गङ्गां (गङ्गाविश्वेश्वरम्) आवाहयामि। ॐ भुवः गङ्गां (गङ्गाविश्वेश्वरम्)
आवाहयामि। ॐ सुवः\ldots। ॐ भूर्भुवस्सुवरोम्, सशक्ति-साङ्ग-सायुध-सवाहन-सर्वालङ्कारभूषितां
श्रीगङ्गां (श्रीगङ्गाविश्वेश्वरम्) आवाहयामि। आवाहिता भव। संस्थापिता भव।\ldots\ सान्निध्यं कुरु।

\twolineshloka*
{ॐ व्योमगामिन्यै विद्महे नारायण्यै च धीमहि}
{तन्नो गङ्गा प्रचोदयात्}
\centerline{\textbf{इदमर्घ्यम् इदमर्घ्यम्।}}

अस्य श्रीगङ्गामूलमन्त्रस्य नारद ऋषिः। गायत्री छन्दः। श्रीगङ्गा देवता। ह्रां इत्यादि
कर-षडङ्गन्यासः। एवं हृदयादिन्यासः।

\dnsub{ध्यानम्}

\fourlineindentedshloka*
{सितमकरनिषण्णां शुभ्रवर्णां त्रिनेत्रां}
{करधृतकलशोद्यत्सोत्पलाभीत्यभीष्टाम्}
{विधिहरिहररूपां सेन्दुकोटीरचूडां}
{भसितसितदुकूलां जाह्नवीं तां नमामि}

\dnsub{मूलमन्त्रः}

\begin{center}
\textbf{ॐ ह्रीं शिवजटाधिरूढायै त्रिपथगामिन्यै गङ्गायै नमः॥}
\end{center}

इति त्रिवारं जप्त्वा, ॐ आं ह्रीं क्रों यं रं लं\ldots\ श्रीगङ्गायाः प्राणा इह प्राणाः।
ॐ आं ह्रीं\ldots\ श्रीगङ्गायाः जीव इह स्थितः। ॐ आं ह्रीं क्रों\ldots\ श्रीगङ्गायाः
सर्वेन्द्रियाणि\ldots\ तिष्ठन्तु स्वाहा। ॐ असुनीते\ldots\ स्वस्ति। ततः पञ्चदशप्रणवावृत्तीः कृत्वा,
ॐ लं पृथिव्यात्मिकायै गन्धं\ldots\ इत्यादि पञ्चोपचारान् कृत्वा।

\sect{द्वादशनाम-पूजा}

\begin{center}
ॐ गङ्गायै नमः। ॐ जह्नुकन्यायै नमः।\\
ॐ भीष्मजनन्यै नमः। ॐ आदिपितामहकमण्डलुकणिकायै नमः।\\
ॐ काशीप्रान्तविहारिण्यै नमः। ॐ शुभ्रवर्णायै नमः।\\
ॐ शिवायै नमः। ॐ भगीरथानुगायै नमः।\\
ॐ तरलतरतरङ्गायै नमः। (ॐ भागीरथ्यै नमः।)\\
ॐ पितृप्रियायै नमः। ॐ मन्दाकिन्यै नमः।\\
ॐ वियन्निम्नगायै नमः।\\
\textbf{श्रीगङ्गायै नमः। द्वादशनामपूजां समर्पयामि॥}
\end{center}

श्रीगङ्गायै नमः, ध्यायामि, ध्यानं समर्पयामि, इत्यादि षोडशोपचारान् कृत्वा।

\sect{कल्पोक्त-षोडशोपचार-पूजा}

{\small प्रत्युपचारं (कोष्ठे) शतरुद्रीयमन्त्रैः श्रीगङ्गाविश्वेश्वरस्यापि।\par}

\renewcommand{\devAya}{श्रीगङ्गायै नमः,}
\begin{center}

\dnsub{ध्यानम्}

\fourlineindentedshloka*
{ब्रह्माण्डं खण्डयन्तीं हरशिरसि जटावल्लिमुल्लासयन्तीं}
{स्वर्लोकादापतन्तीं कनकगिरिगुहागण्डशैलात् स्खलन्तीम्}
{क्षोणीपृष्ठे लुठन्तीं दुरितचयचमूं निर्भरं भर्त्सयन्तीं}
{पाथोधिं पूरयन्तीं सुरनगरसरित्पावनीं भावयामि}

\fourlineindentedshloka*
{आदावादिपितामहस्य निगमव्यापारपात्रे जलं}
{पश्चात् पन्नगशायिनो भगवतः पादोदकं पावनम्}
{भूयः शम्भुजटाविभूषणमणिर्जह्नोर्महर्षेरियं}
{कन्या कल्मषनाशिनी भगवती भागीरथी पातु माम्}
\textbf{\devAya{} ध्यानं समर्पयामि॥}

\dnsub{आवाहनम्}

\fourlineindentedshloka*
{क्वचिन्मुनिगणस्तुते क्वचिदनन्तसम्पूजिते}
{क्वचित्कलकलस्वने क्वचिदधीरयादोगणे}
{क्वचिद्रविकरोज्ज्वले क्वचिदुदग्रपाताकुले}
{भगीरथनृपस्तुते शिवकरं तवावाहनम्}

\twolineshloka*
{विष्णुपादाब्जसम्भूते विश्वनाथशिरःस्थिते}
{आवाहयामि गङ्गे त्वां भक्ताभीष्टफलप्रदे}
ॐ हिरण्यवर्णां\ldots\ आवह॥\\
\textbf{\devAya{} आवाहनं समर्पयामि॥}\\
(ॐ नमस्ते रुद्र\ldots\ ते नमः। श्रीगङ्गाविश्वेश्वराय नमः, आवाहनं समर्पयामि।)

\dnsub{आसनम्}

\twolineshloka*
{गङ्गावारि मनोहारि मुरारिचरणच्युतम्}
{त्रिपुरारिशिरश्चारि पापहारि पुनातु माम्}

\twolineshloka*
{मुक्तारत्नसुवर्णादिखचितं सुन्दरं शुभम्}
{सिंहासनं प्रदास्यामि गृहाण मकरासने}
ॐ तां म आवह\ldots\ पुरुषानहम्॥\\
\textbf{\devAya{} आसनं समर्पयामि॥}\\
(ॐ या त इषुः\ldots\ मृडय। श्रीगङ्गाविश्वेश्वराय नमः, आसनं समर्पयामि।)

\dnsub{पाद्यम्}

\fourlineindentedshloka*
{चरद्विहगसङ्कुले सगरवंशमुक्तिप्रदे}
{मुनीन्द्रवरसंस्तुते दिवि च माल्यसंशोभिते}
{सदा दुरितनाशिके विरचितं त्वदर्थे मया}
{गृहाण जगदम्बिके सुविमलेऽत्र पाद्यं मुदा}

\twolineshloka*
{सिन्ध्वादिसरिदुद्भूतं गन्धपुष्पसमन्वितम्}
{पाद्यं ददाम्यहं देवि प्रसीद परमेश्वरि}
ॐ अश्वपूर्वां\ldots\ जुषताम्॥\\
\textbf{\devAya{} पाद्यं समर्पयामि॥}\\
(ॐ या ते रुद्र\ldots\ अभिचाकशीहि। श्रीगङ्गाविश्वेश्वराय नमः, पाद्यं समर्पयामि।)

\dnsub{अर्घ्यम्}

\twolineshloka*
{पद्मासनस्थे देवेशि काश्यामुत्तरवाहिनि}
{गृहाणार्घ्यं मया दत्तं जाह्नव्यै ते नमो नमः}
ॐ कां सोस्मितां\ldots\ श्रियम्॥\\
\textbf{\devAya{} अर्घ्यं समर्पयामि॥}\\
(ॐ यामिषुं\ldots\ जगत्। श्रीगङ्गाविश्वेश्वराय नमः, अर्घ्यं समर्पयामि।)

\dnsub{आचमनीयम्}

\twolineshloka*
{नमः शिवायै गङ्गायै शिवदायै नमो नमः}
{नमस्ते रुद्ररूपिण्यै गृहाणाचमनीयकम्}

\twolineshloka*
{सुवर्णकलशानीतं नानागन्धसुवासितम्}
{आचम्यतां मया दत्तं गृहाणामृतवर्षिणि}
ॐ चन्द्रां प्रभासां\ldots\ त्वां वृणे॥\\
\textbf{\devAya{} आचमनीयं समर्पयामि॥}\\
(ॐ शिवेन वचसा\ldots\ असत्। श्रीगङ्गाविश्वेश्वराय नमः, आचमनीयं समर्पयामि।)

\dnsub{मधुपर्कः}

\twolineshloka*
{दधिक्षीरघृतैर्युक्तं स्वादूदकसमन्वितम्}
{मधुपर्कं मया दत्तं गृहाण परमेश्वरि}
\textbf{\devAya{} मधुपर्कं समर्पयामि॥}\\
(ॐ अध्यवोचत्\ldots\ यातुधान्यः। श्रीगङ्गाविश्वेश्वराय नमः, मधुपर्कं समर्पयामि।)\\
(ॐ आपो हि ष्ठा\ldots\ चनः। मलापकर्षणस्नानं समर्पयामि।)\\
(अथ `आप्यायस्व' इत्यादिभिः पञ्चामृताभिषेकं प्रोक्षणं वा कुर्यात्।)

\twolineshloka*
{क्षीरं दधि घृतं चैव मधुशर्करयान्वितम्}
{पञ्चामृतं गृहाणेदं प्रसीद परमेश्वरि}

\dnsub{स्नानम् (महाभिषेको वा)}

\fourlineindentedshloka*
{सकलकलुषभङ्गे स्वर्गसोपानसङ्गे}
{तरलतरतरङ्गे देवि गङ्गे प्रसीद}
{विनिहितमिदमत्र स्नानवारि प्रसन्नं}
{विधिहरिहररूपे स्नाहि सौख्यं महेशि}

\twolineshloka*
{नर्मदायमुनासिन्धुगोदावर्याहृतैर्जलैः}
{स्नापयामि शिवे भक्त्या भागीरथि नमोऽस्तु ते}
ॐ आदित्यवर्णे\ldots\ अलक्ष्मीः॥\\
{\small महाभिषेकस्नानं (श्रीरुद्र-चमक-श्रीसूक्त-दुर्गासूक्त-पुरुषसूक्तादिमन्त्रैः यथासम्भवं स्नपयेत् प्रोक्षयेद्वा।)}\\
\textbf{\devAya{} स्नानं समर्पयामि॥}\\
(ॐ असौ यस्ताम्रः\ldots\ ईमहे। श्रीगङ्गाविश्वेश्वराय नमः, स्नानं समर्पयामि।)\\
स्नानानन्तरं पुनराचमनं समर्पयामि॥

\dnsub{वस्त्रम् (वस्त्रयुग्मम्)}

\fourlineindentedshloka*
{सितमकरनिषण्णे भीष्ममातासि नूनं}
{चकितमिव गभीरं देवि रूपं सुधीरम्}
{तव परमहिमानं देवि वाचाममानं}
{भसितसितदुकूले (सुजाती) जाह्नवि त्वं गृहाण}

\twolineshloka*
{वैडूर्यपद्मरागादिखचितं मेखलान्वितम्}
{सुवर्णसूत्रसंयुक्तं वस्त्रं दास्यामि गृह्यताम्}
ॐ उपैतु मां\ldots\ ददातु मे॥\\
\textbf{\devAya{} वस्त्रं (वस्त्रयुग्मं) समर्पयामि॥}\\
(ॐ असौ योऽवसर्पति\ldots\ मृडयाति नः। श्रीगङ्गाविश्वेश्वराय नमः, वस्त्रयुग्मं समर्पयामि।)

\dnsub{उपवीतम्}

\twolineshloka*
{ब्रह्मविष्णुमहेश्वरपाणिप्रोततन्तुरचितं शुभरूपम्}
{यज्ञसूत्रमिदमत्र तवार्थे कल्पितं करुणयाऽनुगृहाण}

\twolineshloka*
{ब्रह्मसूत्रमिदं देवि त्रिगुणीकृतमुज्ज्वलम्}
{उपवीतं गृहाण त्वं विधिशक्रादिपूजिते}
\textbf{\devAya{} उपवीतं समर्पयामि॥}\\
(ॐ नमो अस्तु नीलग्रीवाय\ldots\ अकरं नमः। श्रीगङ्गाविश्वेश्वराय नमः, उपवीतं समर्पयामि।)

\dnsub{आभरणम्}

\fourlineindentedshloka*
{नमो जह्नुकन्ये न मन्ये त्वदन्यां}
{वसिष्ठादिभिर्गीयमानाभिधेयाम्}
{चलन्नूपुरं चाङ्गदद्वन्द्वमेतत्}
{त्वदर्थे प्रदत्तं गृहाणाभिरामम्}

\fourlineindentedshloka*
{माणिक्यमुक्ताफलविद्रुमैश्च}
{गोमेदवैडूर्यकपुष्यरागैः}
{प्रवालनीलैश्च कृतं गृहाण}
{दिव्यं हि रत्नाभरणं च देवि}
ॐ क्षुत्पिपासामलां\ldots\ मे गृहात्॥\\
\textbf{\devAya{} आभरणं समर्पयामि॥}\\
(ॐ नमः शिवाय। श्रीगङ्गाविश्वेश्वराय नमः, आभरणं समर्पयामि।)

\dnsub{गन्धः}

\twolineshloka*
{मलयाचलसम्भूतं श्रीगन्धं चन्द्रमिश्रितम्}
{सुगन्धोज्ज्वलकस्तूरीकुङ्कुमाङ्कितमर्पये}

\twolineshloka*
{मलयाचलसम्भूतं कस्तूरीकुङ्कुमान्वितम्}
{कर्पूरगन्धसम्मिश्रं गृहाण परमेश्वरि}
ॐ गन्धद्वारां\ldots\ ह्वये श्रियम्॥\\
\textbf{\devAya{} गन्धं समर्पयामि॥}\\
(ॐ प्रमुञ्च\ldots\ वप। श्रीगङ्गाविश्वेश्वराय नमः, गन्धं समर्पयामि।)

श्रीकुङ्कुमदिग्धाङ्ग्यै नमः, कुङ्कुमं समर्पयामि।\\
श्रीकस्तूरीतिलकाञ्चितायै नमः, कस्तूरिकां समर्पयामि।\\
नानापरिमलद्रव्यं समर्पयामि।\\
श्रीसौभाग्यदायिन्यै नमः, सौभाग्यद्रव्यं समर्पयामि। रोचनं समर्पयामि।\\
श्रीहरिद्रारञ्जितायै नमः, हरिद्रां समर्पयामि।\\
श्रीसीमन्तभूषणायै नमः, सिन्दूरं समर्पयामि।\\
घनसारं समर्पयामि। कृष्णागरुं समर्पयामि। कक्कोलं समर्पयामि।\\
पट्टसूत्रं समर्पयामि। चूडामणिं समर्पयामि। नासाभरणं समर्पयामि।\\
कङ्कतिकां समर्पयामि। काञ्चीं समर्पयामि। मञ्जीरं समर्पयामि।\\
ताटङ्कं समर्पयामि। कटकं समर्पयामि। यक्षकर्दमं समर्पयामि।\\
अङ्गरागं समर्पयामि। नानाभूषणानि समर्पयामि।\\
\textbf{\devAya{} मङ्गलद्रव्याणि समर्पयामि॥}

\dnsub{अक्षताः}

\twolineshloka*
{अक्षतान् शालिसम्भूतान् हरिद्राकुङ्कुमान्वितान्}
{गृहाण जह्नुतनये अक्षय्यफलदायिनि}
ॐ मनसः काममाकूतिं\ldots\ श्रयतां यशः॥\\
\textbf{\devAya{} अक्षतान् समर्पयामि॥}\\
(ॐ नमः शिवाय। श्रीगङ्गाविश्वेश्वराय नमः, अक्षतान् समर्पयामि।)

\dnsub{पुष्पम्}

\twolineshloka*
{कल्हारोत्पलचाम्पेयजातीमन्दारकेतकैः}
{अर्हणां ते करिष्यामि पुनीहि मम पातकम्}
\textbf{\devAya{} अर्हणापुष्पं समर्पयामि॥}

\twolineshloka*
{केतकीतुलसीबिल्वमल्लिकाकमलादिभिः}
{पुन्नागैरर्चयामि त्वां गृहाणामरवन्दिते}
ॐ कर्दमेन\ldots\ मालिनीम्॥\\
\textbf{\devAya{} पुष्पाणि समर्पयामि॥}

\end{center}

\sect{अङ्ग-पूजा}

\begin{center}
\begin{supertabular}{ll}
ॐ पापपर्वतनाशिन्यै नमः & पादौ पूजयामि\\
ॐ भक्तवत्सलायै नमः & गुल्फौ पूजयामि\\
ॐ जगद्धात्र्यै नमः & जङ्घे पूजयामि\\
ॐ जाह्नव्यै नमः & जानुनी पूजयामि\\
ॐ शैलसुतायै नमः & ऊरू पूजयामि\\
ॐ समुद्रगामिन्यै नमः & कटिं पूजयामि\\
ॐ मकरारूढायै नमः & गुह्यं पूजयामि\\
ॐ आनन्दवर्धिन्यै नमः & जघनं पूजयामि\\
ॐ गङ्गायै नमः & नाभिं पूजयामि\\
ॐ जगत्कुक्ष्यै नमः & उदरं पूजयामि\\
ॐ विशालवक्षायै नमः & वक्षःस्थलं पूजयामि\\
ॐ ह्लादिन्यै नमः & हृदयं पूजयामि\\
ॐ सुस्तन्यै नमः & स्तनौ पूजयामि\\
ॐ तरङ्गिण्यै नमः & पार्श्वे पूजयामि\\
ॐ उन्नतकण्ठ्यै नमः & कण्ठं पूजयामि\\
ॐ त्रैलोक्यसुन्दर्यै नमः & स्कन्धौ पूजयामि\\
ॐ अमृतकलशहस्तायै नमः & हस्तान् पूजयामि\\
ॐ लीलाशुकधारिण्यै नमः & बाहून् पूजयामि\\
ॐ विद्याप्रकाशिन्यै नमः & मुखं पूजयामि\\
ॐ त्रैलोक्यगामिन्यै नमः & ललाटं पूजयामि\\
ॐ सुनासिकायै नमः & नासिकां पूजयामि\\
ॐ मकरकुण्डलधारिण्यै नमः & श्रोत्रे पूजयामि\\
ॐ बिम्बोष्ठ्यै नमः & ओष्ठौ पूजयामि\\
ॐ अनाथरक्षिण्यै नमः & अधरं पूजयामि\\
ॐ चञ्चलगत्यै नमः & जिह्वां पूजयामि\\
ॐ अलकनन्दायै नमः & गण्डस्थलं पूजयामि\\
ॐ तिलकधारिण्यै नमः & फालं पूजयामि\\
ॐ ज्ञानरूपिण्यै नमः & चुबुकं पूजयामि\\
ॐ अमृतबिम्बायै नमः & अलकान् पूजयामि\\
ॐ भीमरथ्यै नमः & शिरः पूजयामि\\
ॐ भागीरथ्यै नमः & सर्वाण्यङ्गानि पूजयामि\\
\end{supertabular}

\textbf{अङ्गपूजां समर्पयामि॥}
\end{center}

\sect{आवरण-पूजा}

\begin{center}
\dnsub{प्रथमावरणम्}
ॐ ह्रां हृदयाङ्गदेवताभ्यो नमः। ॐ ह्रीं शिरोऽङ्गदेवताभ्यो नमः।\\
ॐ हूं शिखाङ्गदेवताभ्यो नमः। ॐ हैं कवचाङ्गदेवताभ्यो नमः।\\
ॐ ह्रौं नेत्राङ्गदेवताभ्यो नमः। ॐ ह्रः अस्त्राङ्गदेवताभ्यो नमः।\\
इति प्रथमावरणम्॥

\dnsub{द्वितीयावरणम्}
ॐ कश्यपाय नमः। ॐ अत्रये नमः। ॐ भरद्वाजाय नमः। ॐ विश्वामित्राय नमः।\\
ॐ गौतमाय नमः। ॐ जमदग्नये नमः। ॐ वसिष्ठाय नमः। ॐ अरुन्धत्यै नमः।\\
इति द्वितीयावरणम्॥

\dnsub{तृतीयावरणम्}
ॐ ब्राह्म्यै नमः। ॐ माहेश्वर्यै नमः। ॐ कौमार्यै नमः। ॐ वैष्णव्यै नमः।\\
ॐ वाराह्यै नमः। ॐ नारसिंह्यै नमः। ॐ इन्द्राण्यै नमः। ॐ चामुण्डायै नमः।\\
इति तृतीयावरणम्॥

\dnsub{चतुर्थावरणम्}
ॐ इन्द्राय नमः। ॐ अग्नये नमः। ॐ यमाय नमः। ॐ निर्ऋतये नमः।\\
ॐ वरुणाय नमः। ॐ वायवे नमः। ॐ सोमाय नमः। ॐ ईशानाय नमः।\\
इति चतुर्थावरणम्॥

\dnsub{पञ्चमावरणम्}
ॐ वज्राय नमः। ॐ शक्तये नमः। ॐ दण्डाय नमः। ॐ खड्गाय नमः।\\
ॐ पाशाय नमः। ॐ अङ्कुशाय नमः। ॐ गदायै नमः। ॐ त्रिशूलाय नमः।\\
इति पञ्चमावरणम्॥

आवरणदेवताभ्यो नमः, सर्वान् जलगन्धाद्युपचारान् समर्पयामि।\\
\textbf{आवरणपूजां समर्पयामि॥}
\end{center}

\sect{पत्र-पूजा}

\begin{center}
ॐ गङ्गायै नमः, पारिजातपत्रं समर्पयामि।\\
ॐ कमलवासिन्यै नमः, मल्लिकापत्रं समर्पयामि।\\
ॐ सर्वज्ञायै नमः, पुन्नागपत्रं समर्पयामि।\\
ॐ वरदायै नमः, अशोकपत्रं समर्पयामि।\\
ॐ भवनाशिन्यै नमः, मरुगपत्रं समर्पयामि।\\
ॐ अभीष्टदायै नमः, दूर्वापत्रं समर्पयामि।\\
ॐ सर्वमङ्गलायै नमः, सर्वाणि पत्राणि समर्पयामि॥
\end{center}

\sect{पुष्प-पूजा}

\begin{center}
ॐ त्रिपथगामिन्यै नमः, शतपत्रपुष्पं समर्पयामि।\\
ॐ भीमरथ्यै नमः, पुन्नागपुष्पं समर्पयामि।\\
ॐ नर्मदायै नमः, अशोकपुष्पं समर्पयामि।\\
ॐ सरस्वत्यै नमः, मल्लिकापुष्पं समर्पयामि।\\
ॐ गोदावर्यै नमः, सेवन्तिकापुष्पं समर्पयामि।\\
ॐ कालिन्द्यै नमः, चम्पकपुष्पं समर्पयामि।\\
ॐ विश्वमात्रे नमः, सर्वाणि पुष्पाणि समर्पयामि॥
\end{center}

\sect{सहस्रनाम-पूजा}

अस्य श्रीगङ्गासहस्रनामस्तोत्रमन्त्रस्य, स्कन्द ऋषिः, अनुष्टुप् छन्दः, श्रीगङ्गाभागीरथी देवता,
मम इष्टार्थसिद्धये पुष्पार्चने विनियोगः॥ ॐ ह्रां ह्रीं इत्यादि षडङ्गन्यासं कुर्यात्॥

\dnsub{ध्यानम्}

\begin{center}
\twolineshloka*
{चतुर्भुजां त्रिनेत्रां च नदीनदनिषेविताम्}
{लावण्यामृतनिष्यन्दसंशीलद्गात्रयष्टिकाम्}
\twolineshloka*
{पूर्णकुम्भस्थिताम्भोजवरदाभयसत्कराम्}
{चामरैर्वीज्यमानां तां श्वेतच्छत्रोपशोभिताम्}
\twolineshloka*
{सुधाप्लावितभूपृष्ठां दिव्यगन्धानुलेपनाम्}
{त्रैलोक्यपूजितां वन्दे देवर्षिभिरभिष्टुताम्}
\end{center}

(श्रीगङ्गादेवीप्रीत्यर्थं सहस्रनामभिः पुष्पार्चने विनियोगः। सहस्रनामार्चने सर्वत्र आदौ `ॐ'
इति योजयित्वा पठ्यते।)

\begingroup
\centering
\setlength{\columnseprule}{1pt}
\let\chapt\sect
\input{../namavali-manjari/1000/Ganga_1000.tex}
\endgroup

\begin{center}
\textbf{इति सहस्रनामपूजां समर्पयामि॥}
\end{center}

\sect{अष्टोत्तरशतनाम-पूजा}

\begingroup
\centering
\setlength{\columnseprule}{1pt}
\let\chapt\sect
% The book's 108 is the list of Ganga_108_2 with a different close (names 95-108);
% which to print here is open (review/ganga/NOTES.md).
\input{../namavali-manjari/100/Ganga_108_2.tex}
\endgroup

\begin{center}
\textbf{इति अष्टोत्तरशतनामपूजां समर्पयामि॥}

\dnsub{धूपः}

\twolineshloka*
{चन्दनागरुमुस्तादिघृतगुग्गुलुसंयुतम्}
{दशाङ्गद्रव्यसंयुक्तं धूपोऽयं प्रतिगृह्यताम्}
कर्दमेन\ldots\ मालिनीम्॥ वनस्पतिरसोत्पन्नः\ldots\ प्रतिगृह्यताम्॥\\
धूरसि धूर्व धूर्वन्तं\ldots\ देवानामसि॥\\
\textbf{\devAya{} धूपम् आघ्रापयामि॥}\\
(ॐ विज्यं धनुः\ldots\ निषङ्गथिः। श्रीगङ्गाविश्वेश्वराय नमः, धूपम् आघ्रापयामि।)

\dnsub{दीपः}

आपः सृजन्तु\ldots\ कुले॥ साज्यं त्रिवर्तिसंयुक्तं\ldots\ तिमिरापहम्।\\
उद्दीप्यस्व\ldots\ दिशो दश॥\\
\textbf{\devAya{} दीपं दर्शयामि॥}\\
(ॐ या ते हेतिः\ldots\ परि ब्भुज। श्रीगङ्गाविश्वेश्वराय नमः, दीपं दर्शयामि।)

\dnsub{नैवेद्यम्}

ॐ आर्द्रां पुष्करिणीं\ldots\ म आवह॥

\twolineshloka*
{शाल्यन्नं व्यञ्जनैर्युक्तं घृतमापसमन्वितम्}
{क्षीरान्नममृतोपेतं भुज्यताममृताशनि}
नैवेद्यं कृत्वा, स्थलशुद्धिं कृत्वा।\\
(ॐ नमस्ते अस्त्वायुधाय\ldots\ धन्वने। श्रीगङ्गाविश्वेश्वराय नमः, नैवेद्यं समर्पयामि।)

\twolineshloka*
{(नरौघाघविनाशिन्यै तरङ्गिण्यै नमो नमः}
{तुभ्यमापोशनं दास्ये स्वीकुरु श्रीहरप्रिये)}
(उत्तरापोशनं समर्पयामि।)\\
हस्तप्रक्षालनं समर्पयामि। मुखप्रक्षालनं समर्पयामि।

\dnsub{ताम्बूलम्}

\twolineshloka*
{पूगीफलसमायुक्तं नागवल्लीदलैर्युतम्}
{कर्पूरचूर्णसंयुक्तं ताम्बूलं प्रतिगृह्यताम्}
ॐ आर्द्रां यः करिणीं\ldots\ म आवह॥\\
\textbf{\devAya{} क्रमुकताम्बूलं समर्पयामि॥}\\
(ॐ परि ते धन्वनः\ldots\ धेहि तम्। श्रीगङ्गाविश्वेश्वराय नमः, ताम्बूलं समर्पयामि।)\\
\textbf{\devAya{} महादक्षिणां समर्पयामि॥ \devAya{} फलनैवेद्यं समर्पयामि॥}

\dnsub{मङ्गलनीराजनम्}

ॐ तां म आवह\ldots\ पुरुषानहम्॥ ॐ हिरण्यपात्रं\ldots\ तेजो ददाति॥\\
ॐ आशास्तेऽयं\ldots\ नमो देवेभ्यः। ॐ ध्रुवं ते राजा\ldots\ ध्रुवम्॥\\
इत्यादि मङ्गलार्तिक्यमन्त्रैः नीराजनं कृत्वा।

\dnsub{मन्त्रपुष्पम्}

\twolineshloka*
{मन्त्रपुष्पं प्रयच्छामि भक्ताभीष्टप्रदायिनि}
{मल्लिकाशोकपुन्नागैर्मरुगैः शतपत्रकैः}
ॐ गणानां त्वा\ldots। ॐ योऽपां पुष्पं वेद\ldots\ आयतनवान् भवति।\\
ॐ यो वेदादौ\ldots\ महेश्वरः॥ ॐ यः शुचिः\ldots\ जपेत्॥\\
\textbf{\devAya{} मन्त्रपुष्पं समर्पयामि॥}\\
(ॐ अघोरेभ्योऽथ\ldots\ रुद्ररूपेभ्यः। श्रीगङ्गाविश्वेश्वराय नमः, मन्त्रपुष्पं समर्पयामि।)

\dnsub{प्रदक्षिण-नमस्कारौ}

\twolineshloka*
{गङ्गे त्रिपथगे दिव्ये जाह्नवि त्रिदिवस्थिते}
{प्रदक्षिणं करोमि त्वां प्रणताघौघनाशिनि}
\twolineshloka*
{नमस्ते जगदाधारे विश्वरूपे सदाशिवे}
{विश्वनाथे महादेवि करुणाकरि शाम्भवि}
\twolineshloka*
{यानि कानि च पापानि जन्मान्तरकृतानि च}
{तानि तानि विनश्यन्ति प्रदक्षिणपदे पदे}
\textbf{\devAya{} प्रदक्षिणं समर्पयामि॥}\\
(ॐ सहस्राणि सहस्रधा\ldots\ मुखा कृधि। श्रीगङ्गाविश्वेश्वराय नमः, प्रदक्षिणं समर्पयामि।)

\twolineshloka*
{पाहि पाहि जगद्वन्द्ये भवतापनिवारिणि}
{नमस्तुभ्यं नमस्तुभ्यं नमस्तुभ्यं सदाशिवे}
\twolineshloka*
{नमस्ते देवदेवेशि नमस्ते करुणाकरे}
{नमस्ते जगदानन्दे नमस्तस्यै नमो नमः}
\textbf{\devAya{} नमस्कारं समर्पयामि॥}

\dnsub{अर्घ्यप्रदानम्}

\twolineshloka*
{जह्नुपुत्रि महादेवि जगद्धात्रि महाबले}
{गन्धपुष्पफलैर्युक्तं गृहाणार्घ्यं नमोऽस्तु ते}
\textbf{\devAya{} इदमर्घ्यम् इदमर्घ्यम् इदमर्घ्यम्॥}\\
एवं त्रिवारम् अर्घ्यं दद्यात्॥

\dnsub{प्रार्थना}

ॐ पद्मानने पद्मऊरु\ldots\ आयुष्मन्तं करोतु माम्॥

\fourlineindentedshloka*
{आदावादिपितामहस्य निगमव्यापारपात्रे जलं}
{पश्चात् पन्नगशायिनो भगवतः पादोदकं पावनम्}
{भूयः शम्भुजटाविभूषणमणिर्जह्नोर्महर्षेरियं}
{कन्या कल्मषनाशिनी भगवती भागीरथी पातु माम्}

\fourlineindentedshloka*
{शैलेन्द्रादवतारिणी निजजले मज्जज्जनोत्तारिणी}
{पारावारविहारिणी भवभयश्रेणीसमुत्सारिणी}
{शेषाहेरनुकारिणी हरशिरोवल्लीदलाकारिणी}
{काशीप्रान्तविहारिणी विजयते गङ्गा मनोहारिणी}

\fourlineindentedshloka*
{मातर्जह्नुमहर्षिपुत्रि वसुधाशृङ्गारहारावलि}
{स्वर्गारोहणवैजयन्ति भवतीं भागीरथीं प्रार्थये}
{त्वत्तीरे वसतस्त्वदम्बु पिबतस्त्वद्वीचिषु प्रेङ्खत-}
{-स्त्वन्नाम स्मरतस्त्वदर्पितदृशः स्यान्मे शरीरव्ययः}

(गङ्गाष्टकैर्वा प्रार्थयेत्।)\\
श्रद्धां मेधां\ldots\ परमेश्वरि॥ मन्त्रहीनं\ldots। अज्ञानाद्वा\ldots॥\\
\textbf{\devAya{} प्रार्थनां समर्पयामि॥}
\end{center}

\sect{उत्तर-(प्रसन्न-)पूजा}

प्रदक्षिण-नमस्कारादि-सर्वोपचारपूजाः समर्पयामि। उत्तर (प्रसन्न) पूजां करिष्ये।\\
ॐ सुमुखाय नमः\ldots\ गणपतये नमः। प्रसन्नपूजां समर्पयामि। प्रसन्नार्घ्यं समर्पयामि॥

\begin{center}
ॐ गङ्गायै नमः। ॐ जह्नुकन्यायै नमः।\\
ॐ भीष्मजनन्यै नमः। ॐ आदिपितामहकमण्डलुकणिकायै नमः।\\
ॐ काशीप्रान्तविहारिण्यै नमः। ॐ शुभ्रवर्णायै नमः।\\
ॐ शिवायै नमः। ॐ भगीरथानुगायै नमः।\\
ॐ तरलतरतरङ्गायै नमः। ॐ पितृप्रियायै नमः।\\
ॐ मन्दाकिन्यै नमः। ॐ वियन्निम्नगायै नमः॥\\
\textbf{ॐ श्रीगङ्गायै नमः। प्रसन्नपूजां (उत्तरपूजां) समर्पयामि। प्रसन्नार्घ्यं समर्पयामि॥}
\end{center}

यस्य स्मृत्या च\ldots\ अच्युतम्॥ कायेन वाचा\ldots\ समर्पयामि।
अस्मिन् वर्तमान\ldots\ श्रीगङ्गाकल्पोक्तपूजाराधनकर्मणा श्रीगङ्गा (श्रीपरमेश्वरः,
श्रीगङ्गाविश्वेश्वरः) प्रीयताम्॥ ॐ तत्सत्। ब्रह्मार्पणमस्तु॥

पूजाकाले स्वरवर्णबिन्दुविसर्ग\ldots\ नामत्रयजपं करिष्ये। ॐ अच्युताय नमः\ldots\
नामत्रयजपं कृत्वा। गोत्रप्रवरमुक्त्वा अभिवाद्य। शङ्खतीर्थं प्रोक्ष्य। `प्रसीद, प्रसादान् देहि'
इति तीर्थप्रसादादिकं सर्वेभ्यो दत्वा, दक्षिणां दत्वा, आशीः सम्प्रार्थ्य, तीर्थप्रोक्षणं प्राशनं
प्रसादग्रहणं च कृत्वा, श्रद्धां मेधां\ldots\ परमेश्वरि\ldots। ॐ उत्तिष्ठ ब्रह्मणस्पते\ldots\ सचा।
यान्तु देवगणाः\ldots\ पुनरागमनाय च। इति घण्टानादपुरस्सरम् उद्वास्य कलशम् आचार्याय दत्वा
आचामेत्॥ यथाशक्ति ब्राह्मणान् सुवासिनीश्च भोजयेत्।

\centerline{॥इति श्रीगङ्गाकल्पोक्तपूजाविधिः॥}
